%% toc: P3 40 40.1
\section{Proofs for probability spaces of finite structures}
\label{sec:annex-p03-w03-p03-064-probability-spaces}

Graphs are irreflexive and symmetric; \(a\not\sim b\) abbreviates
\(\neg\,a\sim b\).  Logarithms are natural.  Two elementary inequalities are
used repeatedly: \(1-x\leq e^{-x}\) for all real \(x\), and
\(\ln(1-x)\geq-x-x^2\) for \(0\leq x\leq1/2\).  The second holds because
\(g(x)=\ln(1-x)+x+x^2\) has \(g(0)=0\) and
\(g'(x)=x(1-2x)/(1-x)\geq0\) on \([0,1/2]\).

\begin{annexitem}{annex:p03-w03-p03-064-yes}{thm:p03-w03-p03-064-yes}{The sentence \(\mathit{YES}\) in \(SU(n,p)\)}
\begin{proof}
\emph{Item 1.}  \(SU(n,p)\not\entails\mathit{YES}\) exactly when no point
satisfies \(U\).  The \(n\) events \(\neg U(x)\) are independent with
probability \(1-p\) each, so this has probability \((1-p)^n\) and
\(\Pr[\mathit{YES}]=1-(1-p)^n\).  As \(p\) increases, \((1-p)^n\) does not
increase; its values at \(p=0\) and \(p=1\) are \(1\) and \(0\).

\emph{Item 2.}  For constant \(p\in(0,1]\), \(0\leq1-p<1\) and
\((1-p)^n\to0\).

\emph{Item 3.}  For \(n>2c\), \(x=c/n\leq1/2\), so
\(-c-c^2/n\leq n\ln(1-c/n)\leq-c\).  Both bounds tend to \(-c\); hence
\((1-c/n)^n\to e^{-c}\) and \(\Pr[\mathit{YES}]\to1-e^{-c}\).

\emph{Item 4.}  By Bernoulli's inequality \((1-p)^n\geq1-np\), so
\(\Pr[\mathit{YES}]\leq np(n)\); if \(p\ll n^{-1}\) then \(np(n)\to0\).  If
\(p\gg n^{-1}\), then \(np(n)\to\infty\) and
\((1-p)^n\leq e^{-np(n)}\to0\).  The two cases are the definition of a
threshold function (\cref{def:p03-w03-p03-065-threshold}) with
\(p_0(n)=n^{-1}\).
\end{proof}
\end{annexitem}

\begin{annexitem}{annex:p03-w03-p03-064-threshold-ratio}{thm:p03-w03-p03-064-threshold-ratio}{Two thresholds for one property}
\begin{proof}
Suppose \(p_1/p_0\to\infty\); the case \(p_1/p_0\to0\) is the same with the
roles exchanged.  Put \(q(n)=\sqrt{p_0(n)p_1(n)}\) for all large \(n\) and
\(q(n)=1/2\) otherwise.  For large \(n\), \(q(n)\in(0,1]\) because \(p_0\) and
\(p_1\) take values in \((0,1]\) there.  Then
\(q/p_0=\sqrt{p_1/p_0}\to\infty\), so \(q\gg p_0\), and since \(p_0\) is a
threshold function, \(A\) holds almost surely at \(q\).  Also
\(q/p_1=\sqrt{p_0/p_1}\to0\), so \(q\ll p_1\), and since \(p_1\) is a
threshold function, \(A\) holds almost never at \(q\).  A sequence of
probabilities cannot tend to both \(1\) and \(0\).
\end{proof}
The hypothesis that the values lie in \((0,1]\) for large \(n\) is what makes
\(q\) an admissible edge or point probability.
\end{annexitem}

\begin{annexitem}{annex:p03-w03-p03-064-unary-boundary}{thm:p03-w03-p03-064-unary-boundary}{The unary zero--one law fails at both thresholds}
\begin{proof}
Item~1 is item~3 of \cref{thm:p03-w03-p03-064-yes}.  For item~2, the events
\(\neg U(x)\) are independent with probability \(1-p(n)=c/n\), so
\(\Pr[\exists x\,\neg U(x)]=1-(1-c/n)^n\), which tends to \(1-e^{-c}\) by the
computation in item~3 of \cref{thm:p03-w03-p03-064-yes}.  Since
\(0<c<\infty\), \(0<1-e^{-c}<1\).
\end{proof}
\end{annexitem}

\begin{annexitem}{annex:p03-w03-p03-064-path3}{thm:p03-w03-p03-064-path3}{Threshold for paths of length three}
Fix \(n\geq10\) and \(p<1\); the case \(p=1\) gives the complete graph, in
which \(P_3\) holds for \(n\geq4\).  The coordinates are the indicators of
the \(\binom n2\) pairs.  For an unordered pair \(a=\{x_1,x_2\}\) of distinct
vertices, with \(x_1\) the smaller label, let \(F_a\) be the set of ordered
pairs \((y_1,y_2)\) of distinct vertices outside \(a\); there are
\(N:=(n-2)(n-3)\) of them.  To \((y_1,y_2)\in F_a\) attach the edge set
\(E=\{x_1y_1,\,y_1y_2,\,y_2x_2\}\) and the event \(B\) that all three edges
are present, of probability \(p^3\).  Distinct members of \(F_a\) have
distinct edge sets, since the edge of \(E\) at \(x_1\) determines \(y_1\) and
the edge at \(x_2\) determines \(y_2\).  Let \(F_a^+\) be \(F_a\) together with
one more index whose edge set is \(\{x_1x_2\}\).  Then \(a\) is joined by a
path of length three exactly when some \(B\) with index in \(F_a\) occurs,
and \(a\) is nonadjacent and joined by no such path exactly when no event
with index in \(F_a^+\) occurs; call the latter event \(\mathrm{Bad}_a\).

\emph{Claim (overlaps).}  Let \(a,b\) be unordered pairs, equal or not, and
let \(i\in F_a^+\), \(j\in F_b^+\) be distinct indices whose edge sets
\(E_i,E_j\) share an edge.  Let \(K\) be the graph formed by
\(E_i\cup E_j\), with \(m\) edges and \(f\) vertices outside \(a\cup b\).
Then \(m\leq5\) and \(f\leq m-2\).  Consequently, writing
\(\Delta(a,b)\) for the sum of \(p^{|E_i\cup E_j|}\) over such ordered pairs
\((i,j)\),
\[
  \Delta(a,b)\leq C_1\bigl(np^3+n^2p^4+n^3p^5\bigr),\qquad C_1:=4\cdot9^4.
\]
\begin{proof}
Each of \(E_i,E_j\) is a path or a single edge, so it spans a connected
graph, and as they share an edge \(K\) is connected; \(K\) contains
\(a\cup b\).  Since \(E_i\neq E_j\) share an edge and have at most three edges
each, \(m\leq5\).  If \(a\neq b\), then \(|a\cup b|\geq3\) and connectedness
gives \(m\geq|V(K)|-1\geq3+f-1\).  If \(a=b\), neither index is the single
edge \(x_1x_2\), because no path of \(F_a\) contains that edge; so \(E_i\) and
\(E_j\) are two different paths from \(x_1\) to \(x_2\).  A graph with two
different paths between two vertices is not a forest, so \(m\geq|V(K)|=2+f\).
In both cases \(f\leq m-2\).

For the count, an ordered pair \((i,j)\) is fixed by its at most four inner
vertices.  Record for each of the four inner slots whether it is absent (an
edge index), equal to one of the at most four vertices of \(a\cup b\), or the
\(k\)-th new vertex, \(k\leq4\); there are at most \(9^4\) such patterns, and
a pattern with \(f\) new vertices is realized by at most \(n^f\) pairs.  So
pairs with given \(m\) contribute at most \(9^4\sum_{f\leq m-2}n^fp^m\leq
9^4(m-1)n^{m-2}p^m\).  Summing over \(m\in\{3,4,5\}\) (with \(m\geq3\)
because \(m\geq f+2\geq2\) and two distinct sets sharing an edge, one of
which has three edges, span at least three) gives the bound.
\end{proof}

\emph{Decreasing events.}  If \(D_1,\ldots,D_k\) are decreasing events in
independent \(\{0,1\}\)-coordinates, then
\(\Pr[D_1\wedge\cdots\wedge D_k]\geq\prod_i\Pr[D_i]\).  For \(k=2\) apply
\cref{thm:p03-w03-p03-065-harris} to the nondecreasing indicator of
\(\neg D_1\) and the nonincreasing indicator of \(D_2\):
\(\Pr[D_2]-\Pr[D_1\wedge D_2]\leq(1-\Pr[D_1])\Pr[D_2]\).  Induct, using that a
conjunction of decreasing events is decreasing.

\begin{proof}
\emph{Item 1.}  Fix \(a\) and apply \cref{thm:p03-w03-p03-065-janson} to the
family \(F_a\): \(\mu=Np^3\), and \(\Delta\leq\Delta(a,a)\).  For
\(n\geq10\), \(N\geq n^2/2\) and \(N\geq(1-5/n)n^2\).

Case (i): \(p\leq n^{-2/3}\ln n\).  Then
\(np^3\leq n^{-1}\ln^3n\), \(n^2p^4\leq n^{-2/3}\ln^4n\) and
\(n^3p^5\leq n^{-1/3}\ln^5n\), so \(\Delta\to0\) uniformly in this case.
By hypothesis \(\mu\geq(1-5/n)(2+\varepsilon)\ln n\geq(2+\varepsilon/2)\ln n\)
for large \(n\).  The first Janson bound gives
\(\Pr[\text{no path for }a]\leq e^{\Delta/2}n^{-2-\varepsilon/2}\).

Case (ii): \(p>n^{-2/3}\ln n\).  Then \(np\geq1\), so \(np^3\leq n^3p^5\) and
\(n^2p^4\leq n^3p^5\), whence \(\Delta\leq3C_1n^3p^5\); and
\(\mu\geq n^2p^3/2\geq\ln^3n/2\).  If \(\Delta<\mu\), the first bound gives
\(e^{-\mu/2}\leq e^{-\ln^3n/4}\).  If \(\Delta\geq\mu\), the second gives
\(e^{-\mu^2/(2\Delta)}\leq\exp\bigl(-np/(24C_1)\bigr)\leq
\exp\bigl(-n^{1/3}\ln n/(24C_1)\bigr)\).

In every case the probability that a fixed pair has no path is at most
\(\max\bigl(e^{\Delta/2}n^{-2-\varepsilon/2},\,e^{-\ln^3n/4},\,
e^{-n^{1/3}\ln n/(24C_1)}\bigr)\), and there are fewer than \(n^2\) pairs.
The union bound gives a failure probability tending to \(0\).  A path of
length three satisfies the displayed sentence, so \(P_3^{\mathrm w}\) holds
whenever \(P_3\) does.

\emph{Item 2.}  Here \(p\leq(2\ln n)^{1/3}n^{-2/3}\to0\); take \(n\) so large
that \(p\leq1/2\).  Then \(np^3\leq2n^{-1}\ln n\),
\(n^2p^4\leq2p\ln n\) and \(n^3p^5=(n^2p^3)(np^2)\leq
(2\ln n)^{5/3}n^{-1/3}\), so \(\delta_n:=4C_1(np^3+n^2p^4+n^3p^5)\to0\).
Let \(X\) be the number of unordered pairs \(a\) for which
\(\mathrm{Bad}_a\) occurs, and \(\mu^+:=p+Np^3\).

\emph{Lower bound for one pair.}  \(\mathrm{Bad}_a\) is the conjunction of
the decreasing events \(\neg B\), so
\(\Pr[\mathrm{Bad}_a]\geq(1-p)(1-p^3)^N\geq
\exp(-\mu^+-p^2-Np^6)\).  As \(Np^6\leq(n^2p^3)^2n^{-2}\to0\) and
\(\mu^+\leq(2-\varepsilon)\ln n+p\),
\(\mathbb E[X]\geq\binom n2e^{-o(1)}n^{-(2-\varepsilon)}\to\infty\).

\emph{Two pairs.}  For \(a\neq b\), apply \cref{thm:p03-w03-p03-065-janson}
to the index set \(F_a^+\sqcup F_b^+\): its \(\mu\) is \(2\mu^+\) and its
\(\Delta\) is at most \(\Delta(a,a)+\Delta(b,b)+2\Delta(a,b)\leq\delta_n\) by
the claim.  So \(\Pr[\mathrm{Bad}_a\wedge\mathrm{Bad}_b]\leq
e^{-2\mu^++\delta_n/2}\), and with the lower bound,
\begin{align*}
  \Pr[\mathrm{Bad}_a\wedge\mathrm{Bad}_b]&\leq(1+\eta_n)\Pr[\mathrm{Bad}_a]
  \Pr[\mathrm{Bad}_b],\\
  1+\eta_n&:=\exp\bigl(\delta_n/2+2p^2+2Np^6\bigr),
\end{align*}
with \(\eta_n\to0\) uniformly in \(a,b\).  By
\cref{thm:p03-w03-p03-065-moments}, applied with \(D\) the set of all ordered
pairs \(a\neq b\),
\(\operatorname{Var}[X]\leq\mathbb E[X]+\eta_n\mathbb E[X]^2\), and
\(\Pr[X=0]\leq1/\mathbb E[X]+\eta_n\to0\).  So almost surely some pair is
nonadjacent and joined by no path of length three.  For a nonadjacent pair,
any \(y_1,y_2\) with \(x_1\sim y_1\sim y_2\sim x_2\) are distinct from each
other and from \(x_1,x_2\): \(y_1\neq x_1\) and \(y_2\neq x_2\) by
irreflexivity, \(y_1\neq y_2\) likewise, and \(y_1=x_2\) or \(y_2=x_1\) would
make \(x_1\sim x_2\).  So such a pair also violates \(P_3^{\mathrm w}\).

\emph{Threshold.}  If \(p\gg n^{-2/3}\ln^{1/3}n\) then
\(n^2p^3/\ln n\to\infty\), so item~1 applies with \(\varepsilon=1\).  If
\(p\ll n^{-2/3}\ln^{1/3}n\) then \(n^2p^3/\ln n\to0\), so item~2 applies with
\(\varepsilon=1\).
\end{proof}
\end{annexitem}

\begin{annexitem}{annex:p03-w03-p03-064-rigid-basic}{thm:p03-w03-p03-064-rigid-basic}{Rigid extensions: monotonicity and enlarging the base}
Throughout, \(e_G(S,Y)\) may be infinite when \(X\) is infinite; the
inequality of \cref{def:p03-w03-p03-064-rigid-extension} counts as satisfied
when it is.  A set \(S\) with \(X\subseteq S\subsetneq X\cup Y\) is determined
by \(Q=S\cap Y\subsetneq Y\), and its nonroot set is \(W=Y\setminus Q\neq
\emptyset\).
\begin{proof}
\emph{Item 1.}  Adding edges leaves every nonroot count unchanged and does
not decrease any edge count, so \(v-e\alpha\) does not increase for any
\(S\).  For finite \(X\), the parameters of \((S,G|_{X\cup Y})\) with
\(X\subseteq S\) are \(|Y\setminus S|\) and \(e_G(S,Y)\); edges between two
vertices of \(X\) are never counted because \(X\subseteq S\).  If \(Y\) forms
an \((R,H)\) extension over \(X\) with \((R,H)\) rigid, then, identifying
roots with \(X\) and nonroots with \(Y\), the induced graph contains every
counted edge of \(H\), so \(G[X;Y]\) is rigid.  Conversely, if \(G[X;Y]\) is
rigid then \(Y\) forms a \(G[X;Y]\) extension over \(X\).

\emph{Item 2.}  Let \(X'\subseteq S'\subsetneq X'\cup Y\) and put
\(S:=X\cup(S'\cap Y)\).  The nonroot sets agree:
\((X'\cup Y)\setminus S'=Y\setminus S'=Y\setminus S=:W\), and \(W\neq\emptyset\)
since \(S'\neq X'\cup Y\).  So \(X\subseteq S\subsetneq X\cup Y\).  Every edge
with an end in \(W\) and the other in \(X\cup Y\) also has its other end in
\(X'\cup Y\), so the edge count over \(X'\) is at least the count over \(X\),
and the inequality for \(S\) gives the one for \(S'\).  The nonroot set of
the pair \((X',Y\setminus X')\) is \(Y\setminus X'\neq\emptyset\), which has at
most \(|Y|\) vertices.

\emph{Item 3 (finite bases).}  For each \(Q\subsetneq Y\), with
\(W=Y\setminus Q\), let \(c_Q\) be the number of edges from \(W\) into
\(X\cup Y\), possibly infinite, and choose a finite \(F_Q\subseteq X\) such
that the number of edges from \(W\) into \(F_Q\cup Y\) is at least
\(\min(c_Q,\lfloor|W|/\alpha\rfloor+1)\).  Let \(X_0\) be the union of the
finitely many \(F_Q\).  For \(X_0\subseteq X'\subseteq X\) and each \(Q\), the
count \(c'_Q\) over \(X'\) satisfies
\(\min(c_Q,\lfloor|W|/\alpha\rfloor+1)\leq c'_Q\leq c_Q\).  If
\(c'_Q>|W|/\alpha\), the inequality holds for \(Q\).  Otherwise
\(c'_Q<\lfloor|W|/\alpha\rfloor+1\), so \(c_Q\) is finite and equals
\(c'_Q\), and the inequality for \(Q\) is the one already given over \(X\).
Thus \(Y\) is rigid over every such \(X'\); in particular over a finite
\(X_0\) containing all neighbours in \(X\) of \(Y\) when there are finitely
many, as stated in the lemma.
\end{proof}
\end{annexitem}

\begin{annexitem}{annex:p03-w03-p03-064-closure-well-defined}{thm:p03-w03-p03-064-closure-well-defined}{The two definitions of closure agree}
A \emph{step sequence} from \(X\) is a finite sequence
\(X=X_0\subseteq X_1\subseteq\cdots\subseteq X_k\) with
\(X_{i+1}=X_i\cup Y_i\), \(Y_i\) disjoint from \(X_i\), \(|Y_i|\leq t\), and
\(Y_i\) rigid over \(X_i\).  A set is \emph{reachable} if it is the last term
of a step sequence; \(C\) denotes the union of all reachable sets.
\begin{proof}
\emph{Steps stay inside closed sets.}  Let \(Z\supseteq X\) be \(t\)-closed
and suppose \(X_i\subseteq Z\).  If \(Y_i\not\subseteq Z\), then by item~2 of
\cref{thm:p03-w03-p03-064-rigid-basic} the set \(Y_i\setminus Z\), of at most
\(t\) vertices, is rigid over \(Z\), contradicting closedness.  So
\(X_{i+1}\subseteq Z\), and by induction every reachable set, hence \(C\),
lies in \(Z\).

\emph{Reachable sets are directed.}  Let \(X_k\) and \(X'_l\) be reachable,
the second through steps \(Y'_0,\ldots,Y'_{l-1}\).  Starting from \(X_k\),
apply in turn the sets \(Y'_j\setminus U_j\), where \(U_j\) is the current
set, skipping \(j\) when this set is empty.  By item~2 of
\cref{thm:p03-w03-p03-064-rigid-basic}, since \(X'_j\subseteq U_j\), each
applied set is rigid over \(U_j\), and it has at most \(t\) vertices.  The
final set is \(X_k\cup X'_l\), which is therefore reachable.

\emph{\(C\) is \(t\)-closed.}  Suppose \(W\) with \(|W|\leq t\), disjoint from
\(C\), is rigid over \(C\).  By item~3 of
\cref{thm:p03-w03-p03-064-rigid-basic} there is a finite \(C_0\subseteq C\)
such that \(W\) is rigid over every \(X'\) with \(C_0\subseteq X'\subseteq C\).
Each vertex of \(C_0\) lies in some reachable set, and by directedness some
reachable \(X_k\) contains \(C_0\).  Then \(W\) is rigid over \(X_k\), so
\(X_k\cup W\) is reachable and \(W\subseteq C\), a contradiction.

\emph{Conclusions.}  \(C\) is \(t\)-closed, contains \(X\), and lies inside
every \(t\)-closed set containing \(X\); so it is the least such set, and it
is defined without reference to any order of steps.  If \(G\) is finite, a
step sequence that cannot be prolonged ends at a reachable \(t\)-closed set,
which lies in \(C\) and, being \(t\)-closed, contains \(C\).  A
\(t'\)-closed set is \(t\)-closed when \(t\leq t'\), so \(cl_{t'}(X)\) is a
\(t\)-closed set containing \(X\) and contains \(cl_t(X)\).  Likewise
\(cl_t(X')\) is \(t\)-closed and contains \(X\) when \(X\subseteq X'\).
\end{proof}
\end{annexitem}

\begin{annexitem}{annex:p03-w03-p03-064-closure-fo}{thm:p03-w03-p03-064-closure-fo}{Closure types are first-order}
Call the finite graph \(H\) with distinguished \(x_1,\ldots,x_r\)
\emph{\(t\)-reachable} if \(V(H)\) is reachable from \(\{x_1,\ldots,x_r\}\)
inside \(H\) by steps of at most \(t\) vertices (for \(t=0\): if
\(V(H)=\{x_1,\ldots,x_r\}\)).
\begin{proof}
If \(H\) is not \(t\)-reachable, let \(\theta_{t,H}\) be \(x_1\neq x_1\).  This
is correct: whether \(Y\) is rigid over \(X\) depends only on the edges
between \(Y\) and \(X\cup Y\), so a step sequence of \(G\) from
\(\{a_1,\ldots,a_r\}\) to a finite \(cl_t(\bar a)\) is also a step sequence
inside the induced graph on \(cl_t(\bar a)\); if that graph is isomorphic to
\(H\) with the roots in place, \(H\) is \(t\)-reachable.

Let \(H\) be \(t\)-reachable, with \(V(H)=\{x_1,\ldots,x_r,y_1,\ldots,
y_m\}\).  Let \(\Delta_H(\bar x,\bar y)\) be the conjunction of
\(u\neq u'\) for distinct listed vertices and of \(u\sim u'\) or
\(\neg\,u\sim u'\) for each pair, as in \(H\).  A \emph{rigid pattern} over
\(H\) is a rigid rooted graph \((V(H),J)\) with \(1\leq w\leq t\) nonroots
\(z_1,\ldots,z_w\), in which \(J\) agrees with \(H\) on \(V(H)\); there are
finitely many.  For each let \(\pi_J(\bar x,\bar y,\bar z)\) say that the
\(z_k\) are pairwise distinct, distinct from all \(x_i,y_j\), and carry
every edge of \(J\) that has an end among the \(z_k\).  Put
\[
  \theta_{t,H}(\bar x):=\exists y_1\cdots\exists y_m\Bigl(\Delta_H(\bar x,
  \bar y)\wedge\bigwedge_J\neg\exists z_1\cdots\exists z_w\,\pi_J\Bigr).
\]
By item~1 of \cref{thm:p03-w03-p03-064-rigid-basic}, for a finite
\(Z\supseteq\{\bar a\}\) the conjunction over \(J\) holds of an enumeration
of \(Z\) exactly when no set of at most \(t\) vertices is rigid over \(Z\),
that is, when \(Z\) is \(t\)-closed.

If \(cl_t(\bar a)\cong H\) with roots in place, take \(\bar b\) to be the
images of \(\bar y\); then \(\Delta_H\) holds and the closure is
\(t\)-closed, so \(\theta_{t,H}(\bar a)\) holds.  Conversely let
\(\bar b\) witness \(\theta_{t,H}(\bar a)\) and \(Z=\{\bar a,\bar b\}\).  The
induced graph on \(Z\) is a copy of \(H\), and a step sequence inside \(H\)
transfers to a step sequence in \(G\), because each rigidity condition
involves only edges inside \(Z\).  So \(Z\) is reachable from
\(\{\bar a\}\), and \(Z\subseteq cl_t(\bar a)\) by
\cref{thm:p03-w03-p03-064-closure-well-defined}.  \(Z\) is \(t\)-closed and
contains \(\bar a\), so \(cl_t(\bar a)\subseteq Z\).  Hence
\(cl_t(\bar a)=Z\cong H\) with \(a_i\) at \(x_i\).  For \(t=0\) the formula is
\(\Delta_H(\bar x)\) alone.
\end{proof}
\end{annexitem}

\begin{annexitem}{annex:p03-w03-p03-064-nonexistence}{thm:p03-w03-p03-064-nonexistence}{Nonexistence Lemma}
\begin{proof}
Fix a graph \(F\) with \(v\leq M\) vertices and \(e\) edges, \(v-e\alpha<0\).
A copy of \(F\) in \(G(n,p)\) is an injective map \(V(F)\to V\) carrying
edges to edges.  There are at most \(n^v\) injective maps, and each carries
the \(e\) edges of \(F\) to \(e\) distinct pairs, all present with
probability \(p^e\).  The expected number of copies is at most
\(n^vp^e=n^{v-e\alpha}\to0\), so by the first moment bound
(\cref{thm:p03-w03-p03-065-moments}) a copy exists with probability tending
to \(0\).  Up to isomorphism there are finitely many such \(F\), and a finite
union of events of probability tending to \(0\) has probability tending to
\(0\).

If \(cl_t(\emptyset)\neq\emptyset\) in a finite graph \(G\), the first step
of a step sequence from \(\emptyset\) is a set \(Y\), \(|Y|\leq t\), rigid over
\(\emptyset\).  Taking \(S=\emptyset\) in the definition, \(|Y|-e\alpha<0\)
for the number \(e\) of edges of \(G|_Y\), so \(G\) contains \(G|_Y\), a graph
with at most \(t\) vertices and \(v-e\alpha<0\).  By the first part with
\(M=t\), this almost surely does not happen.
\end{proof}
\end{annexitem}

\begin{annexitem}{annex:p03-w03-p03-064-finite-closure}{thm:p03-w03-p03-064-finite-closure}{Finite Closure Lemma}
\begin{proof}
For fixed \(v\), \((e\alpha-v)/v\) increases with \(e\), so over the \(e\)
with \(v-e\alpha<0\) it is least at the least integer \(e_v>v/\alpha\).
Hence \(\epsilon=\min_{1\leq v\leq t}(e_v\alpha-v)/v\) exists and is positive.
Put \(K=\lceil r/\epsilon\rceil\geq1\) and \(M=r+K+t\).  By the Nonexistence
Lemma, almost surely \(G\) contains no subgraph with at most \(M\) vertices and
\(v-e\alpha<0\).  We show that every graph \(G\) with this property satisfies
\(|cl_t(x_1,\ldots,x_r)|<r+K\) for all \(x_1,\ldots,x_r\).

Suppose not, and let \(X=\{x_1,\ldots,x_r\}\), so \(|X|\leq r\).  By
\cref{thm:p03-w03-p03-064-closure-well-defined} a maximal step sequence
\(X=R_0\subsetneq R_1\subsetneq\cdots\) ends at \(cl_t(X)\).  Let \(l\) be
least with \(|R_l|\geq r+K\); then \(l\geq1\) and
\(|R_l|<r+K+t\), because \(|R_{l-1}|<r+K\) and a step adds at most \(t\)
vertices.  Let \(F=G|_{R_l}\) with \(V\) vertices and \(E\) edges.  For
\(1\leq i\leq l\) let \(v_i=|R_i\setminus R_{i-1}|\) and \(e_i\) the number
of edges of \(G\) with an end in \(R_i\setminus R_{i-1}\) and the other in
\(R_i\).  Rigidity of the \(i\)-th step, with \(S=R_{i-1}\), gives
\(v_i-e_i\alpha<0\) with \(1\leq v_i\leq t\); by the choice of \(\epsilon\),
\(v_i-e_i\alpha\leq-\epsilon v_i\).  An edge of \(F\) is counted in \(e_i\)
for at most one \(i\), namely the step at which its later end arrives; so
\(E\geq\sum_ie_i\).  With \(V=|X|+\sum_iv_i\),
\[
  V-E\alpha\leq|X|+\sum_{i=1}^l(v_i-e_i\alpha)\leq|X|-\epsilon(V-|X|)
  \leq r-\epsilon K\leq0,
\]
using \(|X|\leq r\) and \(V-|X|\geq r+K-|X|\geq K\).  Since \(\alpha\) is
irrational and \(V\geq1\), \(V-E\alpha\neq0\), so \(V-E\alpha<0\).  But
\(F\) has fewer than \(M\) vertices, which contradicts the property of \(G\).
\end{proof}
Spencer's write-up allows fewer than \(t\) nonroots per step and ends with
\(r-K\epsilon<0\) \cite[p.~9]{Spencer2007LogicRandomStructuresAuthor}; steps
of exactly \(t\) nonroots occur, and \(r-K\epsilon\) may be \(0\), so the
strict inequality is taken from irrationality.
\end{annexitem}

\begin{annexitem}{annex:p03-w03-p03-064-new-vertices}{thm:p03-w03-p03-064-new-vertices}{Spoiler may be required to pick new vertices}
\begin{proof}
A winning strategy in \(\mathrm{EHR}[G,H;t]\) wins the variant, since the
variant only restricts Spoiler.  Conversely let \(\sigma\) win the variant.
In \(\mathrm{EHR}[G,H;t]\), when Spoiler selects a vertex already selected in
that graph in round \(i\), Duplicator selects the round-\(i\) vertex of the
other graph; otherwise she plays \(\sigma\) against the subsequence of rounds
in which Spoiler selected new vertices.  That subsequence is a play of the
variant of length \(k\leq t\).  Because both graphs have at least \(t\)
vertices, Spoiler could continue it to \(t\) rounds with new vertices, and
\(\sigma\) would win; the winning condition for a play implies it for every
prefix, so the \(k\) pairs of the subsequence satisfy it.  The pairs of the
full play are these pairs with repetitions, so the full play satisfies it.

If in a winning play Spoiler selects in round \(j\) a vertex \(x_j\) different
from every earlier \(x_i\) and Duplicator answers \(y_j=y_i\) for some
\(i<j\), then \(x_j\neq x_i\) while \(y_j=y_i\), and the condition fails.  So
new vertices are answered by new vertices.
\end{proof}
\end{annexitem}

\begin{annexitem}{annex:p03-w03-p03-064-ee-game}{thm:p03-w03-p03-064-ee-game}{Elementary equivalence by finite games}
\begin{proof}
Every first-order sentence has some finite quantifier depth.  So \(G\) and
\(H\) satisfy the same sentences exactly when, for every \(t\), they satisfy
the same sentences of quantifier depth at most \(t\), which by
\cref{thm:p03-ef} holds exactly when Duplicator wins
\(\mathrm{EHR}[G,H;t]\).
\end{proof}
\end{annexitem}

\begin{annexitem}{annex:p03-w03-p03-064-infinite-game}{thm:p03-w03-p03-064-infinite-game}{Finite games and the infinite game differ}
\begin{proof}
\emph{Item 1.}  List \(V(G)=\{g_1,g_2,\ldots\}\) and
\(V(H)=\{h_1,h_2,\ldots\}\), with repetitions if a set is finite.  Let
Spoiler select \(g_k\) in round \(2k-1\) and \(h_k\) in round \(2k\), and let
Duplicator follow her winning strategy.  The set of pairs \((x_i,y_i)\)
respects equality in both directions, so it is the graph of an injective
function \(\Phi\); it is defined at every \(g_k\) and takes every value
\(h_k\); and it preserves adjacency and nonadjacency.  So \(\Phi\) is an
isomorphism.

\emph{Item 2.}  Isomorphisms carry connected components to connected
components.  Every component of \(G\) is finite, while \(H\) has the infinite
component \(K_\omega\); so \(G\not\cong H\), and by item~1 Duplicator has no
winning strategy in the \(\omega\)-round game.

Fix \(t\).  Call a component \emph{small} if it is \(K_m\) with \(m<t\), and
\emph{big} otherwise.  The small components of \(G\) and of \(H\) are the same
graphs \(K_1,\ldots,K_{t-1}\); identify them.  Both graphs have infinitely
many big components, each with at least \(t\) vertices.  Duplicator keeps a
partial bijection \(\beta\) between the big components of \(G\) and of \(H\)
that contain selected vertices.  When Spoiler selects a vertex already
selected, she answers its partner.  When he selects a new vertex \(u\) in a
small component, she selects the same vertex in the other graph.  When he
selects a new \(u\) in a big component \(C\): if \(C\) is in the domain or
range of \(\beta\), she selects an unselected vertex of the corresponding
component, which exists because fewer than \(t\) vertices have been selected;
otherwise she selects a vertex of a big component of the other graph that
contains no selected vertex, and adds the pair of components to \(\beta\).

At the end, equality is respected, since new vertices are answered by new
vertices and repeats by partners.  Two distinct vertices of a disjoint union
of complete graphs are adjacent exactly when they lie in the same component;
corresponding components are the same small component or are paired by the
bijection \(\beta\), and no small component corresponds to a big one.  So
adjacency is respected, and Duplicator wins \(\mathrm{EHR}[G,H;t]\).  By
\cref{thm:p03-w03-p03-064-ee-game}, \(G\) and \(H\) are elementarily
equivalent.
\end{proof}
\end{annexitem}

\begin{annexitem}{annex:p03-w03-p03-064-complete-game}{thm:p03-w03-p03-064-complete-game}{Completeness by games on countable models}
\begin{proof}
If \(T\) is complete, any two models satisfy the same sentences, and
\cref{thm:p03-w03-p03-064-ee-game} gives the games.  Suppose \(T\) is not
complete, and choose a sentence \(A\) such that neither \(A\) nor \(\neg A\)
follows from \(T\).  Then \(T\cup\{A\}\) and \(T\cup\{\neg A\}\) have models
(\cref{thm:fol-completeness}).  The language is finite, so by
\cref{thm:downward-ls} each has a countable model; these models are infinite
because \(T\) has no finite models.  Call them \(G\models A\) and
\(H\models\neg A\).  A quantifier-free sentence in a relational language
with equality and no constants is a Boolean combination of \(\top\) and
\(\bot\), so it or its negation is valid; hence \(A\) has quantifier depth
\(t\geq1\).  As \(G\) and \(H\) disagree on \(A\), Duplicator does not win
\(\mathrm{EHR}[G,H;t]\) by \cref{thm:p03-ef}.  The same argument applies to any
finite relational language, with the game as described in the body.
\end{proof}
\end{annexitem}

\begin{annexitem}{annex:p03-w03-p03-064-alice-complete}{thm:p03-w03-p03-064-alice-complete}{The Alice's Restaurant theory is complete}
\begin{proof}
\emph{No finite models.}  A graph with exactly \(m\geq1\) vertices fails
\(A_{m,0}\): for its \(m\) distinct vertices there is no further vertex.

\emph{Staying alive.}  Let \(G,G'\) be models of \(T\).  After some rounds,
let \(a_1,\ldots,a_q\) be the distinct vertices selected in \(G\) and
\(b_1,\ldots,b_q\) their partners in \(G'\); suppose inductively that
\(a_i\mapsto b_i\) is a bijection preserving adjacency and nonadjacency.  If
Spoiler selects an already selected vertex, Duplicator answers its partner.
If he selects a new \(u\in V(G)\), let \(I=\{i:u\sim a_i\}\) and
\(J=\{1,\ldots,q\}\setminus I\).  The axiom \(A_{|I|,|J|}\) of \(G'\), applied
to the distinct \(b_i\) with \(i\in I\) and \(b_j\) with \(j\in J\), gives
\(z\) distinct from all \(b_i\), adjacent to \(b_i\) for \(i\in I\) and
nonadjacent to \(b_j\) for \(j\in J\); Duplicator selects \(z\).  The
extended map is again a bijection preserving adjacency and nonadjacency.
Moves in \(G'\) are answered symmetrically.  So after \(t\) rounds the
condition holds, and Duplicator wins \(\mathrm{EHR}[G,G';t]\), in particular
on countable models.

\emph{Completeness and the limit law.}  \(T\) is consistent: it has a
countable model by \cref{thm:p03-w03-p03-067-categorical}.  By
\cref{thm:p03-w03-p03-064-complete-game}, \(T\) is complete.  Let \(B\) be a
graph sentence.  If \(T\vdash B\), then \(B\) follows from irreflexivity,
symmetry and finitely many \(A_{r,s}\); the first two hold in every
\(G(n,p)\), and each \(A_{r,s}\) holds almost surely by
\cref{thm:p03-w03-p03-067-witness}, so their conjunction, and hence \(B\),
holds almost surely.  Otherwise \(T\vdash\neg B\) and \(B\) holds almost
never.
\end{proof}
\end{annexitem}

\begin{annexitem}{annex:p03-w03-p03-064-f-laws}{thm:p03-w03-p03-064-f-laws}{Laws for \(v-e\alpha\)}
\begin{proof}
\emph{Item 1.}  The vertices outside \(R\) are those of \(S\setminus R\) and
those of \(V(H)\setminus S\).  An edge of \(H\) with an end outside \(R\)
either has both ends in \(S\), and is then an edge of \(H|_S\) with an end
outside \(R\), or has an end outside \(S\); the two cases exclude each other.
So \(e_H(R)=e_{H|_S}(R)+e_H(S)\), and the identity follows.

\emph{Item 2.}  \(|V(H)\setminus S|\geq1\).  If \(e_H(S)=0\), then
\(f(S,H)=|V(H)\setminus S|>0\).  Otherwise \(f(S,H)=0\) would give
\(\alpha=|V(H)\setminus S|/e_H(S)\), a rational number.

\emph{Item 3.}  There are finitely many \(S\), and each \(f(R,H|_S)\) is
positive by safety, so \(\kappa>0\).  By item~1,
\(f(S,H)=f(R,H)-f(R,H|_S)\leq f(R,H)-\kappa\).

\emph{Item 4.}  As \((R,H)\) is not safe, \(f(R,H|_S)\leq0\) for some
\(R\subsetneq S\subseteq V(H)\), and by item~2 applied to \(H|_S\) the value
is not \(0\).  Choose \(S\) with \(R\subsetneq S\subseteq V(H)\) for which
\(m:=f(R,H|_S)\) is least; then \(m<0\).  Let \(R\subseteq S'\subsetneq S\).
By item~1 applied to \(H|_S\), whose restriction to \(S'\) is \(H|_{S'}\),
\(f(S',H|_S)=m-f(R,H|_{S'})\).  If \(S'=R\), the last term is \(0\) and
\(f(S',H|_S)=m<0\).  Otherwise \(f(R,H|_{S'})\geq m\) by the choice of
\(S\), so \(f(S',H|_S)\leq0\), and by item~2 it is not \(0\).  Hence every
nailextension \((S',H|_S)\) is dense, and \((R,H|_S)\) is rigid.
\end{proof}
\end{annexitem}

\begin{annexitem}{annex:p03-w03-p03-064-moment-bound}{thm:p03-w03-p03-064-moment-bound}{An upper bound for counts of arbitrary extensions}
An \((R,J)\) extension over \(\bar x\) is identified with an injective map
\(\iota\) from \(V(J)\) to the vertex set that sends the roots to
\(x_1,\ldots,x_r\) in order, the nonroots outside \(\{x_1,\ldots,x_r\}\), and
every edge of \(J\) with an end outside \(R\) to an edge of \(G\).  For such
a map let \(E_\iota\) be the set of images of the edges of \(J\) with an end
outside \(R\).

\emph{Structures.}  For a \(k\)-tuple \((\iota_1,\ldots,\iota_k)\) of
injective maps of this kind, its \emph{structure} is the equivalence
relation on \(\{1,\ldots,k\}\times(V(J)\setminus R)\) that relates \((j,a)\)
and \((j',a')\) when \(\iota_j(a)=\iota_{j'}(a')\).  A set of \(m\geq1\)
elements carries at most \(m^m\) equivalence relations, since each is
determined by the map sending an element to the least member of its class.
So there are at most \((kw)^{kw}\) structures.  The structure determines,
for each \(j\), the set \(S_j\) of \(a\in V(J)\) with \(\iota_j(a)\) among the
roots or among the values of \(\iota_1,\ldots,\iota_{j-1}\); it contains
\(R\).  The number of classes is \(c=\sum_j|V(J)\setminus S_j|\), and at most
\(n^c\) tuples have a given structure.

\begin{proof}
Let \(N=N_J(\bar x)\).  Then \(\mathbb E[N^k]\) is the sum, over
\(k\)-tuples of maps as above, of the probability that all pairs in
\(E_{\iota_1}\cup\cdots\cup E_{\iota_k}\) are edges, which is
\(p^{|E_{\iota_1}\cup\cdots\cup E_{\iota_k}|}\) because distinct pairs are
independent.  For each \(j\), the images under \(\iota_j\) of the
\(e_J(S_j)\) edges of \(J\) with an end outside \(S_j\) are distinct pairs
with an end at a vertex that is not a value of any earlier map, so they lie
in no \(E_{\iota_l}\) with \(l<j\); a pair counted at step \(j\) has both ends
among the roots and the values of \(\iota_1,\ldots,\iota_j\), so it is not counted again at
a later step.  Hence \(|\bigcup_jE_{\iota_j}|\geq\sum_je_J(S_j)\), and the
contribution of the tuples of one structure is at most
\[
  \prod_{j=1}^kn^{|V(J)\setminus S_j|}p^{e_J(S_j)}=\prod_{j=1}^kn^{f(S_j,J)}
  \leq n^{kF^+},
\]
since \(f(S_j,J)\leq F^+\) when \(S_j\neq V(J)\) and \(f(V(J),J)=0\leq F^+\).
So \(\mathbb E[N^k]\leq(kw)^{kw}n^{kF^+}\), and by Markov's inequality
\(\Pr[N\geq n^{F^++\delta}]=\Pr[N^k\geq n^{k(F^++\delta)}]\leq
(kw)^{kw}n^{-k\delta}\).  The constant \(C=(kw)^{kw}\) does not depend on
\(n\) or \(\bar x\).
\end{proof}
\end{annexitem}

\begin{annexitem}{annex:p03-w03-p03-064-large-deviation}{thm:p03-w03-p03-064-large-deviation}{Large deviation and counting for safe extensions}
Use the notation of the previous proof with \(J=H\) and \(w=v\).  Let
\(\Omega\) be the set of maps \(\iota\) there described, without the edge
condition; \(|\Omega|=(n-r)(n-r-1)\cdots(n-r-v+1)\).  For \(\iota\in\Omega\)
let \(I_\iota\) be the indicator that every pair of \(E_\iota\) is an edge,
so that \(N=N(\bar x)=\sum_{\iota\in\Omega}I_\iota\),
\(\mathbb E[I_\iota]=p^e\) and \(\mu=|\Omega|p^e\).  Put \(f:=f(R,H)\) and
let \(\kappa>0\) be as in item~3 of \cref{thm:p03-w03-p03-064-f-laws}.
\begin{proof}
\emph{The mean.}  For \(n\geq2(r+v)\) each of the \(v\) factors of
\(|\Omega|\) is at least \(n/2\), so \(\mu\geq2^{-v}n^vp^e=2^{-v}n^f\).

\emph{A central moment.}  Fix \(k\geq1\).  Expanding,
\[
  \mathbb E\bigl[(N-\mu)^{2k}\bigr]=\sum_{(\iota_1,\ldots,\iota_{2k})\in
  \Omega^{2k}}\mathbb E\Bigl[\prod_{l=1}^{2k}(I_{\iota_l}-p^e)\Bigr].
\]
For a tuple, let \(\Gamma\) be the graph on \(\{1,\ldots,2k\}\) in which
\(l\neq l'\) are adjacent when \(E_{\iota_l}\cap E_{\iota_{l'}}\neq\emptyset\).
If \(l\) is isolated in \(\Gamma\), the factor \(I_{\iota_l}-p^e\) is a
function of the edge indicators in \(E_{\iota_l}\), the product of the other
factors is a function of indicators outside \(E_{\iota_l}\), the two are
independent, and the first has mean \(0\); so the term vanishes.  For the
remaining tuples,
\[
  \Bigl|\mathbb E\prod_l(I_{\iota_l}-p^e)\Bigr|\leq
  \mathbb E\prod_l(I_{\iota_l}+p^e)=\sum_{A\subseteq\{1,\ldots,2k\}}
  p^{|\bigcup_{l\in A}E_{\iota_l}|}\,p^{e(2k-|A|)}.
\]
The structure of the tuple, defined as before, determines \(\Gamma\) and the
right-hand side.  Fix a structure in which \(\Gamma\) has no isolated vertex,
and fix \(A\).  Reorder the indices so that the components of \(\Gamma\) come
one after another and, inside each component, every index after the first is
adjacent in \(\Gamma\) to an earlier one; the bound below concerns this order,
and \(S_j\) is computed in it.  As in the previous proof the edges of \(H\)
with an end outside \(S_j\) give \(e_H(S_j)\) pairs new at step \(j\), so
\(|\bigcup_{l\in A}E_{\iota_l}|\geq\sum_{j\in A}e_H(S_j)\); for \(j\notin A\),
\(p^e\leq p^{e_H(S_j)}\) because \(e_H(S_j)\leq e_H(R)=e\).  With at most
\(n^{\sum_j|V(H)\setminus S_j|}\) tuples, the contribution of the structure
and \(A\) is at most \(\prod_jn^{f(S_j,H)}\).

If \(j\) is the first index of its component, \(f(S_j,H)\leq f\): this is
equality for \(S_j=R\) and item~3 of \cref{thm:p03-w03-p03-064-f-laws}
otherwise.  If \(j\) is a later index, \(E_{\iota_j}\) shares a pair with an
earlier \(E_{\iota_l}\); every pair in these sets has an end outside
\(\{x_1,\ldots,x_r\}\), and that end is a value of \(\iota_j\) at a nonroot
and of the earlier map, so \(S_j\neq R\) and \(f(S_j,H)\leq f-\kappa\).  Every
component has at least two elements, so there are \(c\leq k\) components, and
\[
  \sum_jf(S_j,H)\leq cf+(2k-c)(f-\kappa)=2kf-(2k-c)\kappa\leq2kf-k\kappa .
\]
There are at most \((2kv)^{2kv}\) structures and \(2^{2k}\) sets \(A\), so
\[
  \mathbb E\bigl[(N-\mu)^{2k}\bigr]\leq2^{2k}(2kv)^{2kv}\,n^{2kf-k\kappa}.
\]

\emph{Item 1.}  By Markov's inequality applied to \((N-\mu)^{2k}\) and the
bound on the mean, for \(n\geq2(r+v)\),
\[
  \Pr\bigl[|N-\mu|>\varepsilon\mu\bigr]\leq
  \frac{\mathbb E[(N-\mu)^{2k}]}{(\varepsilon\mu)^{2k}}\leq
  2^{2k}(2kv)^{2kv}2^{2kv}\varepsilon^{-2k}\,n^{-k\kappa}.
\]
Choose \(k>r/\kappa\).  The bound is \(o(n^{-r})\), and it does not depend on
\(\bar x\).

\emph{Item 2.}  There are at most \(n^r\) tuples \(\bar x\) of distinct
vertices.  By the uniform bound of item~1, the probability that some
\(\bar x\) has \(|N(\bar x)-\mu|>\varepsilon\mu\) is at most a constant times
\(n^{r-k\kappa}\), which tends to \(0\).
\end{proof}
If \(e=0\), every \(I_\iota\) equals \(1\), \(N=\mu\), and the statements are
trivial; the proof covers this case, since then no tuple has an edge in
\(\Gamma\).
\end{annexitem}

\begin{annexitem}{annex:p03-w03-p03-064-exp-deviation}{thm:p03-w03-p03-064-exp-deviation}{The deviation probability is at most \(\exp[-n^{\lambda}]\)}
Use the notation of the previous proof.
\begin{proof}
\emph{The moment bound for every order.}  The bound
\(\mathbb E[(N-\mu)^{2k}]\leq2^{2k}(2kv)^{2kv}n^{2kf-k\kappa}\) of the
previous proof was derived for an arbitrary integer \(k\geq1\) and an
arbitrary \(n\): the count of structures, the count of sets \(A\), the count
of tuples of one structure and the exponent estimate use no relation between
\(k\) and \(n\).  If \(e=0\), then \(N=\mu\) and there is nothing to prove;
let \(e\geq1\).  For \(n\geq2(r+v)\) the mean satisfies
\(\mu\geq2^{-v}n^f\), so Markov's inequality applied to \((N-\mu)^{2k}\)
gives, for every \(k\geq1\),
\[
  \Pr\bigl[|N-\mu|>\varepsilon\mu\bigr]
  \leq\frac{2^{2k}(2kv)^{2kv}n^{2kf-k\kappa}}{\varepsilon^{2k}2^{-2kv}n^{2kf}}
  =\bigl(C\,k^{2v}\,n^{-\kappa}\bigr)^k,
  \qquad C:=4\,(4v)^{2v}\varepsilon^{-2},
\]
because \((2kv)^{2kv}2^{2kv}=\bigl((4v)^{2v}k^{2v}\bigr)^k\).

\emph{Choice of the order.}  Let \(k=\lfloor n^{\kappa/(4v)}\rfloor\), which
is at least \(1\) for \(n\geq1\).  Then \(k^{2v}\leq n^{\kappa/2}\), so
\(Ck^{2v}n^{-\kappa}\leq Cn^{-\kappa/2}\).  Take \(n_0\geq2(r+v)\) so large
that for \(n\geq n_0\) both \(Cn^{-\kappa/2}\leq e^{-1}\) and
\(n^{\kappa/(4v)}-1\geq n^{\kappa/(8v)}\); the second holds for large \(n\)
because \(n^{\kappa/(4v)}-n^{\kappa/(8v)}=n^{\kappa/(8v)}(n^{\kappa/(8v)}-1)\)
tends to infinity.  For \(n\geq n_0\),
\[
  \Pr\bigl[|N-\mu|>\varepsilon\mu\bigr]\leq e^{-k}\leq
  e^{-(n^{\kappa/(4v)}-1)}\leq\exp[-n^{\kappa/(8v)}]=\exp[-n^{\lambda}],
\]
and the bound does not depend on \(\bar x\).  There are at most \(n^r\)
tuples of distinct roots, and the union bound gives the last statement.
\end{proof}
\end{annexitem}

\begin{annexitem}{annex:p03-w03-p03-064-theory-as}{thm:p03-w03-p03-064-theory-as}{Every sentence of \(T_\alpha\) holds almost surely}
\begin{proof}
\emph{Nonexistence.}  For a graph \(H\) with \(v\) vertices, \(e\) edges and
\(v-e\alpha<0\), apply \cref{thm:p03-w03-p03-064-nonexistence} with
\(M=v\).

\emph{Generic Extension.}  Fix a safe \((R,H)\) with parameters \(r,v,e\)
and \(t\geq0\), and put \(f:=f(R,H)\), which is positive (take
\(S=V(H)\) in the definition of safe).  Let \(\mathcal J\) be the finite set
of rooted graphs \((R,J)\) of the following two kinds.
\begin{enumerate}
\item[(a)] \(V(J)=V(H)\) and \(E(J)=E(H)\cup\{g\}\), where \(g\notin E(H)\)
  joins a nonroot of \(H\) to another vertex of \(H\).
\item[(b)] (Only if \(t\geq1\).)  \(V(J)=V(H)\cup Z\) with \(Z\) a set of
  \(w\) new vertices, \(1\leq w\leq t\); \(E(J)=E(H)\cup E_Z\), where every
  pair in \(E_Z\) has an end in \(Z\); the rooted graph \((V(H),J)\) is
  rigid; and some pair in \(E_Z\) joins \(Z\) to \(V(H)\setminus R\).
\end{enumerate}

\emph{Claim A: failures are counted by \(\mathcal J\).}  Let
\(y_1,\ldots,y_v\) be an \((R,H)\) extension over \(\bar x\).  If it fails
condition~2, some edge of \(G\) joins \(y_i\) to \(y_j\) or to \(x_j\) without
being required by \(H\); adding the corresponding pair to \(H\) gives \(J\)
of kind (a), and \(\bar y\) is an \((R,J)\) extension.  If it fails
condition~3, then \(t\geq1\) and there are distinct \(z_1,\ldots,z_w\),
\(w\leq t\), distinct from the \(x_i\) and \(y_j\), forming a rigid extension
over \(\bar x,\bar y\), with an edge to some \(y_j\).  By item~1 of
\cref{thm:p03-w03-p03-064-rigid-basic}, \(\{z_1,\ldots,z_w\}\) is rigid over
\(\{\bar x,\bar y\}\).  Let \(J\) have \(E_Z\) the set of pairs corresponding
to the edges of \(G\) with an end among the \(z_k\) and the other among the
\(x_i,y_j,z_k\).  Then \((V(H),J)\) has the counted edges of
\(G[\{\bar x,\bar y\};\{\bar z\}]\), so it is rigid; it has a pair joining
\(Z\) to a nonroot; and \((\bar y,\bar z)\) is an \((R,J)\) extension over
\(\bar x\).  Distinct \(\bar y\) give distinct tuples.  Hence the number of
\((R,H)\) extensions over \(\bar x\) failing condition~2 or~3 is at most
\(\sum_{J\in\mathcal J}N_J(\bar x)\).

\emph{Claim B: \(F^+<f\) for every \(J\in\mathcal J\).}  As \(f>0\), it
suffices that \(f(S,J)<f\) for \(R\subseteq S\subsetneq V(J)\).

Kind (a).  Here \(e_J(S)\geq e_H(S)\), so \(f(S,J)\leq f(S,H)\), which is
less than \(f\) for \(S\neq R\) by item~3 of
\cref{thm:p03-w03-p03-064-f-laws}.  For \(S=R\), \(g\) has an end outside
\(R\), so \(f(R,J)=f-\alpha<f\).

Kind (b).  Write \(S_H=S\cap V(H)\).  If \(Z\not\subseteq S\), put
\(S'=S\cup V(H)\), so \(V(H)\subseteq S'\subsetneq V(J)\) and, by item~1 of
\cref{thm:p03-w03-p03-064-f-laws},
\(f(S,J)=f(S,J|_{S'})+f(S',J)\).  The second term is negative because
\((V(H),J)\) is rigid.  In the first, \(S'\setminus S=V(H)\setminus S_H\),
and every edge of \(H\) with an end outside \(S_H\) is an edge of
\(J|_{S'}\) with an end outside \(S\); so \(f(S,J|_{S'})\leq f(S_H,H)\leq f\)
(equality at \(S_H=R\), item~3 for \(R\subsetneq S_H\subsetneq V(H)\), and
\(0\) at \(S_H=V(H)\)).  Thus \(f(S,J)<f\).  If \(Z\subseteq S\), then
\(S_H\neq V(H)\), the vertices outside \(S\) are those of
\(V(H)\setminus S_H\), and \(e_J(S)\) is \(e_H(S_H)\) plus the number of
pairs of \(E_Z\) meeting \(V(H)\setminus S_H\).  If \(S_H\neq R\), then
\(f(S,J)\leq f(S_H,H)<f\).  If \(S_H=R\), some pair of \(E_Z\) meets
\(V(H)\setminus R\), and \(f(S,J)\leq f-\alpha<f\).

\emph{Estimate.}  Let \(\delta>0\) be half of \(f-\max_{J\in\mathcal J}F^+_J\)
(or \(f/2\) if \(\mathcal J\) is empty), and choose an integer \(k\) with
\(k\delta>r\).  For distinct \(\bar x\) consider the events
\(N(\bar x)<\mu/2\) and, for \(J\in\mathcal J\),
\(N_J(\bar x)\geq n^{F^+_J+\delta}\).  By the proof of
\cref{thm:p03-w03-p03-064-large-deviation} with \(\varepsilon=1/2\), and by
\cref{thm:p03-w03-p03-064-moment-bound}, each has probability at most a
constant times \(n^{-r-\eta}\) for some \(\eta>0\), uniformly in \(\bar x\).
There are at most \(n^r\) tuples \(\bar x\) and finitely many \(J\), so
almost surely none of these events occurs.  Then for every \(\bar x\), by
Claims A and B, fewer than \(|\mathcal J|\,n^{f-\delta}\) extensions fail
condition~2 or~3, while \(N(\bar x)\geq\mu/2\geq2^{-v-1}n^f\).  For large
\(n\), \(|\mathcal J|\,n^{f-\delta}<2^{-v-1}n^f\), so some extension satisfies
conditions 1--3.  This proves the Generic Extension sentence almost surely,
and also the stated bound \(O(n^{f-\delta})\) on the failures.
\end{proof}
\end{annexitem}

\begin{annexitem}{annex:p03-w03-p03-064-closure-models}{thm:p03-w03-p03-064-closure-models}{Closures in models of the Nonexistence schema}
Step sequences and reachable sets are those of the proof of
\cref{thm:p03-w03-p03-064-closure-well-defined}, which also shows that
\(cl_t(X)\) is the union of the reachable sets and that the union of two
reachable sets is reachable.
\begin{proof}
\emph{Item 1.}  Let \(X\) be finite and \(X=R_0\subsetneq R_1\subsetneq
\cdots\subsetneq R_l\) a step sequence with \(l\geq1\).  For each \(i\) let
\(v_i=|R_i\setminus R_{i-1}|\), so \(1\leq v_i\leq t\), and let \(e_i\) be the
number of edges of \(G\) with an end in \(R_i\setminus R_{i-1}\) and the other
in \(R_i\).  Rigidity of the step, with \(S=R_{i-1}\), gives
\(v_i-e_i\alpha<0\), so by the definition of \(\epsilon_t\),
\(v_i-e_i\alpha\leq-\epsilon_tv_i\).  Let \(F=G|_{R_l}\), with \(V\)
vertices and \(E\) edges; each edge is counted in at most one \(e_i\), so
\(E\geq\sum_ie_i\) and
\[
  V-E\alpha\leq|X|+\sum_i(v_i-e_i\alpha)\leq|X|-\epsilon_t(V-|X|).
\]
By hypothesis \(V-E\alpha\geq0\), and \(V\geq1\); if \(E=0\) then
\(V-E\alpha=V>0\), and otherwise \(V-E\alpha\neq0\) because \(\alpha\) is
irrational.  So \(\epsilon_t(V-|X|)<|X|\), that is,
\(|R_l\setminus X|<|X|/\epsilon_t\).  For \(X=\emptyset\) this is impossible,
so no step starts at \(\emptyset\) and \(cl_t(\emptyset)=\emptyset\).  For
finite \(X\), every reachable set has fewer than \(|X|+|X|/\epsilon_t\)
elements.  Let \(M\) be a reachable set of greatest size.  For every
reachable \(R'\), \(M\cup R'\) is reachable and no larger than \(M\), so
\(R'\subseteq M\).  Hence \(cl_t(X)=M\), which is finite, reachable and
satisfies the bound.

\emph{Item 2.}  Let \(A=V(F)\) and \(t=|A|\).  The rooted graph
\((\emptyset,G|_A)\) has \(f(\emptyset,G|_A)\leq|V(F)|-|E(F)|\alpha<0\), so it
is not safe.  By item~4 of \cref{thm:p03-w03-p03-064-f-laws} there is a
nonempty \(S\subseteq A\) with \((\emptyset,G|_S)\) rigid; that is, \(S\), of
at most \(t\) vertices, is rigid over \(\emptyset\), and
\(S\subseteq cl_t(\emptyset)\).

\emph{Locality.}  For \(t=0\) both closures are \(X\).  Let \(t\geq1\).
Whether a finite \(Y\) is rigid over a finite \(U\) depends only on the
induced graph on \(U\cup Y\), which is the same in \(G\) and in \(G|_W\) when
\(U\cup Y\subseteq W\).  Every step sequence of \(G\) from \(X\) stays inside
the \(t\)-closed set \(W\) (first paragraph of the proof of
\cref{thm:p03-w03-p03-064-closure-well-defined}).  So the step sequences from
\(X\) in \(G\) and in \(G|_W\) are the same, the reachable sets are the same,
and so are their unions, the two closures.
\end{proof}
\end{annexitem}

\begin{annexitem}{annex:p03-w03-p03-064-duplicator}{thm:p03-w03-p03-064-duplicator}{Duplicator wins on models of \(T_\alpha\)}
\emph{Preliminaries.}  (P1) A model of \(T_\alpha\) contains no copy of any
finite graph with \(v-e\alpha<0\), so
\cref{thm:p03-w03-p03-064-closure-models} applies to it: for \(t\geq1\) and
finite \(X\), \(cl_t(X)\) is finite and reachable from \(X\), with
\(|cl_t(X)\setminus X|<|X|/\epsilon_t\), and \(cl_t(\emptyset)=\emptyset\);
and \(cl_0(X)=X\).  (P2) If \(\psi\) is an isomorphism of an induced subgraph
\(G|_A\) onto \(G'|_{A'}\), then for finite disjoint \(U,Y\subseteq A\), \(Y\)
is rigid over \(U\) exactly when \(\psi(Y)\) is rigid over \(\psi(U)\), since
rigidity depends only on the induced graph on \(U\cup Y\); so \(\psi\) carries
step sequences inside \(A\) to step sequences inside \(A'\).  (P3) If
\(t\leq t'\), a \(t'\)-closed set is \(t\)-closed.

\emph{The parameters.}  The \(t_i\) of the body satisfy
\(0=t_0\leq t_1\leq\cdots\leq t_u\), \(t_i\geq1\) for \(i\geq1\), and for
\(1\leq i\leq u\) and every set \(X^+\) of at most \(u\) vertices of a model
of \(T_\alpha\),
\(|cl_{t_{i-1}}(X^+)\setminus X^+|\leq t_i-1\).  For \(i=1\) the left side is
\(0\).  For \(i\geq2\), by (P1) it is less than
\(u/\epsilon_{t_{i-1}}\leq t_i-1\).

\begin{proof}
For \(s\leq u\) write \(X_s=\{x_1,\ldots,x_s\}\) and
\(X_s'=\{x_1',\ldots,x_s'\}\).  The invariant \((I_s)\) is: with \(i=u-s\),
there is an isomorphism \(\varphi\) of \(G|_{cl_{t_i}(X_s)}\) onto
\(G'|_{cl_{t_i}(X_s')}\) with \(\varphi(x_j)=x_j'\) for \(j\leq s\).

\emph{Start and end.}  \((I_0)\) holds with the empty map, because
\(cl_{t_u}(\emptyset)=\emptyset\) in both graphs by (P1).  If \((I_u)\) holds,
then \(t_0=0\) and \(\varphi\) is an isomorphism of \(G|_{X_u}\) onto
\(G'|_{X_u'}\) with \(x_j\mapsto x_j'\); as a function it gives
\(x_j=x_l\Rightarrow x_j'=x_l'\), as an injective map the converse, and as an
isomorphism \(x_j\sim x_l\iff x_j'\sim x_l'\).  That is Duplicator's win.

\emph{The step.}  Assume \((I_s)\) with \(s<u\), so \(i=u-s\geq1\).  Put
\(b=t_i=\mathit{BIG}\), \(a=t_{i-1}=\mathit{SMALL}\),
\(\mathit{OLD}=cl_b(X_s)\) and \(\mathit{OLD}'=cl_b(X_s')=\varphi(\mathit{OLD})\);
both are finite by (P1).  The invariant is symmetric under exchanging \(G\)
and \(G'\) and replacing \(\varphi\) by \(\varphi^{-1}\), so assume Spoiler
selects \(y\in V(G)\).  Let \(X^+=X_s\cup\{y\}\), of at most \(u\) vertices.

\emph{Inside: \(y\in\mathit{OLD}\).}  Duplicator selects
\(y'=\varphi(y)\); put \(X'^+=X_s'\cup\{y'\}\).  By (P3), \(\mathit{OLD}\) is
\(a\)-closed and contains \(X^+\), and \(\mathit{OLD}'\) is \(a\)-closed and
contains \(X'^+\).  By the locality part of
\cref{thm:p03-w03-p03-064-closure-models}, \(cl_a(X^+)\) and
\(cl_a(X'^+)\) may be computed in \(G|_{\mathit{OLD}}\) and
\(G'|_{\mathit{OLD}'}\).  By (P2), applied to \(\varphi\) and
\(\varphi^{-1}\), \(\varphi\) maps the step sequences from \(X^+\) inside
\(\mathit{OLD}\) onto those from \(X'^+\) inside \(\mathit{OLD}'\), hence
maps the union of the reachable sets onto the union of the reachable sets:
\(\varphi(cl_a(X^+))=cl_a(X'^+)\).  The restriction of \(\varphi\) gives
\((I_{s+1})\).

\emph{Outside: \(y\notin\mathit{OLD}\).}  Let \(C=cl_a(X^+)\), finite by
(P1), and \(\mathit{NEW}=C\setminus\mathit{OLD}\), which contains \(y\).
Since \(X_s\subseteq\mathit{OLD}\),
\(\mathit{NEW}\subseteq(C\setminus X^+)\cup\{y\}\), so
\(|\mathit{NEW}|\leq(b-1)+1=b\) by the choice of the parameters.

(i) \emph{Safety.}  Let \(H=G|_{\mathit{OLD}\cup\mathit{NEW}}\).  If the
rooted graph \((\mathit{OLD},H)\) were not safe, item~4 of
\cref{thm:p03-w03-p03-064-f-laws} would give
\(\mathit{OLD}\subsetneq S\subseteq\mathit{OLD}\cup\mathit{NEW}\) with
\((\mathit{OLD},H|_S)\) rigid; then \(S\setminus\mathit{OLD}\), a nonempty set
of at most \(b\) vertices, would be rigid over \(\mathit{OLD}\), against the
\(b\)-closedness of \(\mathit{OLD}=cl_b(X_s)\).  So \((\mathit{OLD},H)\) is
safe.

(ii) \emph{The answer.}  List \(\mathit{OLD}=\{o_1,\ldots,o_q\}\) and
\(\mathit{NEW}=\{w_1,\ldots,w_m\}\), and regard \((\mathit{OLD},H)\) as a
rooted graph with roots \(o_1,\ldots,o_q\) and nonroots \(w_1,\ldots,w_m\).
Apply in \(G'\) the Generic Extension sentence for it with \(t=a\) at the
distinct vertices \(\varphi(o_1),\ldots,\varphi(o_q)\), which list
\(\mathit{OLD}'\).  It gives distinct \(w_1',\ldots,w_m'\) outside
\(\mathit{OLD}'\) such that, by conditions 1 and 2, \(w_k'\sim w_l'\iff
w_k\sim w_l\) and \(w_k'\sim\varphi(o_j)\iff w_k\sim o_j\) for all \(k,l,j\),
and, by condition~3, \(\mathit{NEW}'=\{w_1',\ldots,w_m'\}\) is \(a\)-generic
over \(\mathit{OLD}'\).  Let \(\psi\) extend \(\varphi\) by \(w_k\mapsto
w_k'\).  Then \(\psi\) is an isomorphism of \(G|_{\mathit{OLD}\cup
\mathit{NEW}}\) onto \(G'|_{D'}\), where \(D'=\mathit{OLD}'\cup\mathit{NEW}'\).
Duplicator selects \(y'=\psi(y)\); put \(X'^+=X_s'\cup\{y'\}=\psi(X^+)\).

(iii) \emph{\(\psi(C)\subseteq cl_a(X'^+)\).}  By (P1), \(C\) is the last set
of a step sequence from \(X^+\); its sets lie in
\(C\subseteq\mathit{OLD}\cup\mathit{NEW}\), so by (P2) their images form a
step sequence from \(X'^+\) ending at \(\psi(C)\).

(iv) \emph{\(\psi(C)\) is \(a\)-closed.}  This is trivial for \(a=0\).  Let
\(a\geq1\), and suppose \(W'\), with \(1\leq|W'|\leq a\) and disjoint from
\(\psi(C)\), is rigid over \(\psi(C)\).  If \(W'\subseteq D'\), then by (P2)
\(\psi^{-1}(W')\) is rigid over \(C\), which contradicts the \(a\)-closedness
of \(C\).  Otherwise \(Z'=W'\setminus D'\) is nonempty, and by item~2 of
\cref{thm:p03-w03-p03-064-rigid-basic} it is rigid over \(D'\).  By item~1
of that lemma its vertices form a rigid extension over the vertices of
\(D'\), and there are at most \(a\) of them, so by \(a\)-genericity no edge
joins \(Z'\) to \(\mathit{NEW}'\).  Then \(Z'\) is rigid over
\(\mathit{OLD}'\): for \(\mathit{OLD}'\subseteq S\subsetneq\mathit{OLD}'\cup
Z'\), the set \(S\cup\mathit{NEW}'\) lies between \(D'\) and \(D'\cup Z'\),
has the same nonroots \(Z'\setminus S\), and, as no edge joins \(Z'\) to
\(\mathit{NEW}'\), the edges counted for it over \(D'\) are exactly those
counted for \(S\) over \(\mathit{OLD}'\); so the rigidity inequality for
\(S\cup\mathit{NEW}'\) is the one for \(S\).  Thus \(Z'\) is a nonempty set of
at most \(a\leq b\) vertices rigid over \(\mathit{OLD}'=cl_b(X_s')\), which
contradicts its \(b\)-closedness by (P3).

(v) \emph{Conclusion.}  By (iv), \(\psi(C)\) is an \(a\)-closed set containing
\(X'^+\), so \(cl_a(X'^+)\subseteq\psi(C)\); with (iii),
\(\psi(C)=cl_a(X'^+)\).  The restriction of \(\psi\) to \(C\) is an
isomorphism of \(G|_C\) onto \(G'|_{cl_a(X'^+)}\) sending \(x_j\) to
\(x_j'\) and \(y\) to \(y'\), which is \((I_{s+1})\).

By induction on \(s\), Duplicator keeps the invariant through all \(u\)
rounds and wins.
\end{proof}
Spencer requires the bound on the number of added vertices for sets of size
\(i+1\) \cite[p.~11]{Spencer2007LogicRandomStructuresAuthor}; the step uses
it for the sets \(X^+\), which have up to \(u\) vertices, and the parameters
above are chosen for that.  The bound \(|\mathit{NEW}|\leq\mathit{BIG}\)
is the one Spencer attributes to the definition of \(\mathit{BIG}\)
\cite[p.~12]{Spencer2007LogicRandomStructuresAuthor}.
\end{annexitem}

\begin{annexitem}{annex:p03-w03-p03-064-sparse-law}{thm:p03-w03-p03-064-sparse-law}{\(T_\alpha\) is complete; the zero--one law}
\begin{proof}
\emph{Consistency.}  A derivation of a contradiction from \(T_\alpha\) would
use finitely many sentences of \(T_\alpha\).  The graph axioms hold in every
\(G(n,p)\), and each of the other finitely many sentences holds almost surely
by \cref{thm:p03-w03-p03-064-theory-as}; so for large \(n\) some graph
satisfies all of them, and by soundness they do not derive a contradiction.

\emph{No finite models.}  For \(m\geq1\) let \(R=\{x_1,\ldots,x_m\}\) and let
\(H\) have one further vertex and no edges.  Then \(f(R,H)=1>0\) and
\(S=V(H)\) is the only set to check, so \((R,H)\) is safe, and its Generic
Extension sentence with \(t=0\) is in \(T_\alpha\).  It says that every \(m\)
distinct vertices are accompanied by a further vertex, which fails in a graph
with exactly \(m\) vertices.

\emph{Completeness.}  By \cref{thm:p03-w03-p03-064-duplicator}, Duplicator
wins \(\mathrm{EHR}[G,G';u]\) for all models, in particular all countable
models, \(G,G'\) of \(T_\alpha\) and all \(u\geq1\).  By
\cref{thm:p03-w03-p03-064-complete-game}, \(T_\alpha\) is complete.

\emph{Zero--one law.}  If \(T_\alpha\vdash A\), then \(A\) is derivable from
finitely many sentences of \(T_\alpha\).  Their conjunction holds almost
surely, since a finite intersection of events with probabilities tending to
\(1\) has probability tending to \(1\), and by soundness \(A\) holds in every
graph satisfying the conjunction; so \(A\) holds almost surely.  Otherwise,
by completeness, \(T_\alpha\vdash\neg A\), so \(\neg A\) holds almost surely
and \(A\) almost never.
\end{proof}
\end{annexitem}

\begin{annexitem}{annex:p03-w03-p03-064-containment}{thm:p03-w03-p03-064-containment}{Threshold for containing a fixed graph}
Let \(H\) have \(v\) vertices and \(e\geq1\) edges, and put \(m=m(H)\).  For
a graph \(K\), let \(X_K\) be the number of copies of \(K\) in
\(G(n,p)\).
\begin{proof}
\emph{Below the threshold.}  Choose a subgraph \(K\) of \(H\) with
\(v_K\geq1\) and \(e_K=mv_K\).  As \(m\geq\frac12\), \(e_K\geq1\).  A copy
of \(H\) restricts to a copy of \(K\), so \(A_H\) implies \(X_K>0\).  There
are at most \(n^{v_K}\) injective maps from \(V(K)\) to \([n]\), each a copy
with probability \(p^{e_K}\), so \(\mathbb E[X_K]\leq n^{v_K}p^{mv_K}
=(np^m)^{v_K}\).  If \(p\ll n^{-1/m}\), then \(p^m\ll n^{-1}\) by
\cref{def:p03-w03-p03-065-asymptotic-order}, that is \(np^m\to0\).  By the
first moment bound of \cref{thm:p03-w03-p03-065-moments},
\(\Pr[A_H]\leq\Pr[X_K>0]\leq\mathbb E[X_K]\to0\).

\emph{Above the threshold.}  Let \(p\gg n^{-1/m}\).  Then \(p>0\) for large
\(n\), and \(1/(np^m)=(n^{-1/m}/p)^m\to0\), so \(np^m\to\infty\).  Write
\(X=X_H=\sum_\varphi I_\varphi\), the sum over injective
\(\varphi:V(H)\to[n]\), where \(E_\varphi\) is the set of images of the
edges of \(H\) and \(I_\varphi\) indicates that every pair of \(E_\varphi\)
is an edge.  For \(n\geq2v\), \(\mathbb E[X]\geq(n/2)^vp^e\geq
2^{-v}(np^m)^v\), since \(e\leq mv\) and \(p\leq1\); so
\(\mathbb E[X]\to\infty\).  Two indicators with
\(E_\varphi\cap E_\psi=\emptyset\) depend on disjoint sets of pairs and are
independent.  By \cref{thm:p03-w03-p03-065-moments}, with \(D\) the set of
ordered pairs \((\varphi,\psi)\), \(\varphi\neq\psi\), with
\(E_\varphi\cap E_\psi\neq\emptyset\),
\[
  \Pr[X=0]\leq\frac{1}{\mathbb E[X]}+\frac{1}{\mathbb E[X]^2}
  \sum_{(\varphi,\psi)\in D}\Pr[I_\varphi=I_\psi=1].
\]
Fix \((\varphi,\psi)\in D\) and let \(U\) be the set of \(x\in V(H)\) with
\(\psi(x)\in\varphi(V(H))\), \(s=|U|\).  A pair in \(E_\varphi\cap E_\psi\)
is \(\psi(\{x,y\})\) for an edge \(xy\) of \(H\), and both of its ends lie in
\(\varphi(V(H))\); so \(x,y\in U\), \(s\geq2\), and
\(|E_\varphi\cap E_\psi|\leq e(H|_U)\leq ms\).  Hence
\(\Pr[I_\varphi=I_\psi=1]=p^{2e-|E_\varphi\cap E_\psi|}\leq p^{2e-ms}\).
For a given \(s\) there are at most \(n^v\) choices of \(\varphi\),
\(2^v\) of \(U\), \(v^v\) of \(\psi\) on \(U\) (into \(\varphi(V(H))\)) and
\(n^{v-s}\) of \(\psi\) on the rest.  With
\(\mathbb E[X]^2\geq4^{-v}n^{2v}p^{2e}\),
\[
  \frac{1}{\mathbb E[X]^2}\sum_{D}\Pr[I_\varphi=I_\psi=1]
  \leq\sum_{s=2}^{v}\frac{2^vv^vn^{2v-s}p^{2e-ms}}{4^{-v}n^{2v}p^{2e}}
  =8^vv^v\sum_{s=2}^{v}(np^m)^{-s}\longrightarrow0 .
\]
So \(\Pr[A_H]=\Pr[X>0]\to1\).
\end{proof}
\end{annexitem}

\begin{annexitem}{annex:p03-w03-p03-064-hab}{thm:p03-w03-p03-064-hab}{The maximum density of \(H_{a,b}\)}
Use the notation of \cref{def:p03-w03-p03-064-hab}, and write \(\mathbf
1_P\) for the indicator of \(P\) on \(\{0,\ldots,v-1\}\).
\begin{proof}
\emph{Counting.}  For \(1\leq j\leq q+1\) we have \(0<2j\leq2q+2<v\).  The
two differences of the pair \(\{i,i+j\}\) modulo \(v\) are \(j\) and
\(v-j>j\), so pairs with different \(j\) are different.  For fixed \(j\),
\(\{i,i+j\}=\{i',i'+j\}\) with \(i\neq i'\) would force \(i=i'+j\) and
\(i'=i+j\), so \(2j\equiv0\pmod v\), which is false.  Telescoping,
\(\sum_{i=0}^{l-1}\mathbf 1_P(i)=\lfloor lr/a\rfloor\) for
\(0\leq l\leq v\); with \(l=v=aM\), \(|P|=rM\), and each run
\(\{ta,\ldots,ta+a-1\}\) contains \(\lfloor(t+1)r\rfloor-\lfloor tr\rfloor=r\)
elements of \(P\).  So \(H_{a,b}\) has \(qv+rM=(qa+r)M=bM\) edges on
\(v=aM\) vertices.

\emph{Weights.}  For \(0\leq i\leq v-1\) let
\(D_i=\sum_{j=1}^{i}\bigl(\mathbf 1_P(j)-r/a\bigr)\), so \(D_0=0\).  Since
\(r<a\), \(\mathbf 1_P(0)=\lfloor r/a\rfloor=0\), and
\(D_i=\lfloor(i+1)r/a\rfloor-(i+1)r/a+r/a\), which lies in
\((r/a-1,\,r/a]\).  Put \(w_i=1-r/a+D_i\in(0,1]\).  Split every edge as two
nonnegative weights with sum \(1\), one at each end:
\begin{itemize}
\item an edge \(\{i,i+j\}\) with \(2\leq j\leq q\): \(\frac12\) to each end;
\item a chord \(\{i,i+q+1\}\) with \(i\in P\): \(1\) to \(i\) and \(0\) to
  \(i+q+1\);
\item an edge \(\{i,i+1\}\), \(0\leq i\leq v-1\): \(w_i\) to \(i+1\) and
  \(1-w_i\) to \(i\).
\end{itemize}
A vertex \(x\) lies in the two edges \(\{x,x+j\}\), \(\{x-j,x\}\) for each
\(2\leq j\leq q\), receiving \(q-1\) in all; it receives \(\mathbf 1_P(x)\)
from chords; and it receives \(1-w_x\) from \(\{x,x+1\}\) and
\(w_{x-1}\) from \(\{x-1,x\}\), indices modulo \(v\).  Its total is
\(q+\mathbf 1_P(x)-(w_x-w_{x-1})\).  For \(1\leq x\leq v-1\),
\(w_x-w_{x-1}=D_x-D_{x-1}=\mathbf 1_P(x)-r/a\).  For \(x=0\),
\(D_{v-1}=\sum_{j=1}^{v-1}\mathbf 1_P(j)-(v-1)r/a
=rM-\mathbf 1_P(0)-(v-1)r/a=r/a-\mathbf 1_P(0)\), so
\(w_0-w_{v-1}=-D_{v-1}=\mathbf 1_P(0)-r/a\).  In every case the total is
\(q+r/a=b/a\).

\emph{The bound.}  For a vertex set \(U\), every edge with both ends in
\(U\) has its two weights at vertices of \(U\), so the number of such edges
is at most the sum over \(x\in U\) of the weight \(x\) receives, which is
\((b/a)|U|\).  A subgraph \(K\) with vertex set \(U\) has at most that many
edges, so \(m(H_{a,b})\leq b/a\); and \(H_{a,b}\) itself has density
\(bM/(aM)=b/a\).
\end{proof}
\end{annexitem}

\begin{annexitem}{annex:p03-w03-p03-064-rational-threshold}{thm:p03-w03-p03-064-rational-threshold}{A threshold at every rational exponent}
\begin{proof}
Since \(q\geq1\), \(H_{a,b}\) has edges, and \(m(H_{a,b})=b/a\) by
\cref{thm:p03-w03-p03-064-hab}.  By \cref{thm:p03-w03-p03-064-containment},
\(n^{-1/m(H_{a,b})}=n^{-a/b}\) is a threshold function for
\(A=A_{H_{a,b}}\).  Let \(\alpha<a/b<\beta\) be irrational.  For
\(p=n^{-\alpha}\), \(p/n^{-a/b}=n^{a/b-\alpha}\to\infty\), so
\(p\gg n^{-a/b}\) and \(A\) holds almost surely.  By
\cref{thm:p03-w03-p03-064-sparse-law}, \(A\) holds almost never unless
\(T_\alpha\vdash A\); hence \(T_\alpha\vdash A\).  For \(p=n^{-\beta}\),
\(p/n^{-a/b}=n^{a/b-\beta}\to0\), so \(A\) holds almost never, \(\neg A\)
almost surely, and by the same corollary \(T_\beta\vdash\neg A\).  As
\(T_\beta\) is consistent, \(T_\beta\nvdash A\).
\end{proof}
\end{annexitem}

\begin{annexitem}{annex:p03-w03-p03-064-rep-layout}{thm:p03-w03-p03-064-rep-layout}{A block layout represents its relation}
\begin{proof}
\emph{Item 1.}  The interval \([c,d)\) consists of the vertices
\(c+j(k+1)+i\) with \(j<m\) and \(0\leq i\leq k\).  In the layout those with
\(i=0\) have empty trace and the others have traces of one element.  Suppose
\(\operatorname{Rep}^k_u(c,d;\bar x)\) holds with witnesses
\(e_0,\ldots,e_k\).  Then \(c\leq e_0<e_1<\cdots<e_k<d\) with
\(e_i=e_0+i\), so all \(e_i\) lie in \([c,d)\), and \(N_u(e_0)=\emptyset\).
Hence \(e_0=c+j(k+1)\) for some \(j<m\), and \(e_i=c+j(k+1)+i\) has trace
\(\{r_{j,i}\}\) for \(1\leq i\leq k\).  The formula says \(x_i\leq u\),
\(e_i\sim x_i\) and that no other \(w\leq u\) is adjacent to \(e_i\), so
\(N_u(e_i)=\{x_i\}\) and \(x_i=r_{j,i}\).  Thus \(\bar x=\bar r_j\in R\).
Conversely, for \(\bar x=\bar r_j\) the vertices \(e_i=c+j(k+1)+i\),
\(0\leq i\leq k\), are consecutive, satisfy \(c\leq e_0\) and
\(e_k\leq c+m(k+1)-1<d\), and have the required traces.

\emph{Item 2.}  With \(c=d=u+1\), no \(e_0\) satisfies
\(c\leq e_0\leq e_k<d\).
\end{proof}
\end{annexitem}

\begin{annexitem}{annex:p03-w03-p03-064-every-relation}{thm:p03-w03-p03-064-every-relation}{Every \(k\)-ary relation is represented almost surely}
Logarithms here are to base \(2\).  \emph{Trace fact}: for vertices
\(e>u\), the trace \(N_u(e)\) is determined by the \(u\) pairs \(\{w,e\}\),
\(w\leq u\); for distinct such \(e\) these sets of pairs are disjoint, so the
traces are independent, and each prescribed trace has probability \(2^{-u}\)
in \(G_<(n,\frac12)\).
\begin{proof}
Write \(U=U_k(n)\) and \(n'=\lfloor n/2\rfloor\).  Fix \(u\) with
\(1\leq u\leq U\) and \(R\subseteq\{1,\ldots,u\}^k\).  If \(R=\emptyset\),
\cref{thm:p03-w03-p03-064-rep-layout}, item~2, gives \(c=d=u+1\) for every
graph, once \(n>u\).  Let \(R\) have \(m\geq1\) elements and put
\(s=m(k+1)\).  Since \(m\leq u^k\leq u^{k+1}\leq\log n/(2k+2)\),
\[
  us\leq(k+1)u^{k+1}\leq\tfrac12\log n,\qquad s\leq(k+1)\log n .
\]
Let \(K=\lfloor(n'-u)/s\rfloor\) and \(c_j=u+1+js\) for \(0\leq j<K\).  The
blocks \([c_j,c_j+s)\) are disjoint subsets of \((u,n']\), and
\(c_j+s\leq n'+1\leq n\).  Let \(W_j\) be the event that a block layout of
\(R\) is at \(c_j\); it prescribes the traces of \(s\) vertices above \(u\),
so by the trace fact \(\Pr[W_j]=2^{-us}\geq n^{-1/2}\), and the \(W_j\) are
independent.  Hence
\[
  \Pr\Bigl[\bigcap_{j<K}\neg W_j\Bigr]\leq(1-n^{-1/2})^K\leq\exp(-Kn^{-1/2}).
\]
For \(n\) at least some \(n_1\) depending only on \(k\),
\(n/4\geq1+u+s\) holds for all \(u\leq U\) (both \(u\) and \(s\) are at most
\((k+1)\log n\)), and then \(K\geq(n/2-1-u)/s-1\geq n/(4s)\geq
n/(4(k+1)\log n)\).  If some \(W_j\) occurs, \(c=c_j\) and \(d=c_j+s\)
represent \(R\) by \cref{thm:p03-w03-p03-064-rep-layout}, item~1.

There are at most \(U\leq\log n\) values of \(u\), and for each at most
\(2^{u^k}\leq2^{\log n}=n\) relations.  So the probability that some pair
\((u,R)\) is not represented is at most
\[
  n\log n\,\exp\Bigl(-\frac{n^{1/2}}{4(k+1)\log n}\Bigr),
\]
which tends to \(0\).
\end{proof}
\end{annexitem}

\begin{annexitem}{annex:p03-w03-p03-064-influence-basics}{thm:p03-w03-p03-064-influence-basics}{Two facts about total influence}
For \(x\in\{0,1\}^V\), \(x^{\oplus b}\) is \(x\) with coordinate \(b\)
flipped.
\begin{proof}
\emph{Item 1.}  If no \(g_i\) has \(b\) pivotal at \(x\), then
\(g_i(x)=g_i(x^{\oplus b})\) for all \(i\), so the AND and the OR of the
\(g_i\) take the same value at \(x\) and \(x^{\oplus b}\).  Hence
\(\{x:b\text{ pivotal for }f\}\) is contained in the union of the sets
\(\{x:b\text{ pivotal for }g_i\}\), and the inequality follows by the union
bound for the uniform \(x\).

\emph{Item 2.}  Negating \(f\) changes no pivotal coordinate, and the
negation of a CNF of width at most \(w\) is a DNF of width at most \(w\)
(De Morgan: exchange AND and OR and negate the literals).  So let \(f\) be a
DNF with terms \(T_1,\ldots,T_m\).  Let \(\mathrm{piv}(x)\) be the number of
coordinates pivotal at \(x\); then
\(\mathbf I[f]=2^{-|V|}\sum_x\mathrm{piv}(x)\).  Each pair
\(\{x,x^{\oplus b}\}\) on which \(f\) changes is counted twice in this sum,
once at each end, and exactly one end has value \(1\); so
\(\sum_x\mathrm{piv}(x)=2\sum_{f(x)=1}\mathrm{piv}(x)\).  If \(f(x)=1\), some
\(T_i\) is satisfied by \(x\); if \(b\) is pivotal at \(x\), then
\(f(x^{\oplus b})=0\), so \(T_i(x^{\oplus b})=0\) and \(b\) is a variable of
\(T_i\).  Hence \(\mathrm{piv}(x)\leq w\) when \(f(x)=1\), and
\(\mathbf I[f]\leq2\cdot2^{-|V|}\cdot w\cdot|f^{-1}(1)|\leq2w\).
\end{proof}
\end{annexitem}

\begin{annexitem}{annex:p03-w03-p03-064-depth-induction}{thm:p03-w03-p03-064-depth-induction}{Total influence of narrow bounded-depth circuits}
\begin{proof}
Induction on \(d\).  For \(d=2\), \(f\) is computed by a DNF or CNF of width
at most \(w\), and \(\mathbf I[f]\leq2w=c_2w\Lambda^0\) by
\cref{thm:p03-w03-p03-064-influence-basics}, item~2.

Let \(d\geq3\) and let the circuit \(C\) compute \(f\).  Each layer-\(2\)
gate \(g\), as a function of the input, is computed by a DNF or CNF of width
at most \(w\) (its inputs are layer-\(1\) gates of the other label); let
\(S_2\leq S\) be the number of these gates.  Put \(\delta=1/(32w)\leq\frac12\)
and \(t=\lceil\log_2(2wS^2)\rceil\geq1\), let \(\rho=(J\mid z)\) be a
\(\delta\)-random restriction, write \(h_\rho\) for \(h_{J\mid z}\), and let \(B\) be the event that some
layer-\(2\) gate has \(\operatorname{DT}(g_\rho)\geq t\).  By
\cref{thm:p03-w03-p03-071-switching} and the union bound,
\(\Pr[B]\leq S_2(16\delta w)^t=S_22^{-t}\leq S/(2wS^2)=1/(2wS)\).

\emph{On \(B\).}  The function \(f_\rho\) is computed from the functions
\(g_\rho\) of the layer-\(2\) gates by the gates of layers \(3\) to \(d\).
Applying item~1 of \cref{thm:p03-w03-p03-064-influence-basics} gate by gate
from the output downwards, a coordinate pivotal for \(f_\rho\) at \(y\) is
pivotal for some \(g_\rho\) at \(y\); so
\(\mathbf I[f_\rho]\leq\sum_g\mathbf I[g_\rho]\).  Substituting the fixed
values into a DNF of width at most \(w\) (deleting falsified terms and true
literals) gives a DNF of width at most \(w\), and dually for a CNF; so
\(\mathbf I[g_\rho]\leq2w\) by item~2 of
\cref{thm:p03-w03-p03-064-influence-basics} and \(\mathbf I[f_\rho]\leq2wS\).

\emph{Off \(B\).}  Suppose the layer-\(2\) gates are OR gates, so the
layer-\(3\) gates are AND gates; the other case is the same with AND and OR
exchanged and DNFs in place of CNFs.  By item~3 of
\cref{thm:p03-w03-p03-071-restriction-facts}, each \(g_\rho\) is computed by
a CNF of width at most \(t-1\); by \cref{def:p03-w03-p03-071-restriction}
its clauses involve distinct variables.  Build a
circuit \(C'\) over the variables of \(J\): layer \(0\) is the \(2|J|\)
literals; layer \(1\) has an OR gate for each clause of each such CNF; layer
\(2\) has, for each layer-\(3\) gate \(h\) of \(C\), an AND gate whose
inputs are all clauses of the CNFs of the inputs of \(h\); and for
\(l\geq3\), layer \(l\) of \(C'\) consists of the gates of layer \(l+1\) of
\(C\) with the same connections.  The AND of the clauses of several CNFs is
the AND of the functions they compute, so the layer-\(2\) gate for \(h\)
computes \(h_\rho\), and by induction on layers \(C'\) computes \(f_\rho\).
The labels alternate, \(C'\) has depth \(d-1\) and width at most \(t\), and
its layers \(2,\ldots,d-2\) are layers \(3,\ldots,d-1\) of \(C\), with at
most \(S\) gates.  By the induction hypothesis,
\(\mathbf I[f_\rho]\leq c_{d-1}t\Lambda'^{\,d-3}\) with
\(\Lambda'=\log_2(4tS)\).  Now \(t\leq\log_2(4wS^2)\leq2\Lambda\); since
\(\Lambda\geq2\), \(\log_2(2\Lambda)\leq\Lambda\); and
\(2+\log_2S\leq\Lambda\).  So
\(\Lambda'=2+\log_2t+\log_2S\leq2+\log_2(2\Lambda)+\log_2S\leq2\Lambda\),
and \(\mathbf I[f_\rho]\leq2^{d-2}c_{d-1}\Lambda^{d-2}\).

\emph{Conclusion.}  By \cref{thm:p03-w03-p03-071-restriction-facts}, item~1,
\begin{align*}
  \mathbf I[f]&=\frac1\delta\mathbb E\bigl[\mathbf I[f_\rho]\bigr]\\
  &\leq32w\Bigl(2^{d-2}c_{d-1}\Lambda^{d-2}+\Pr[B]\cdot2wS\Bigr)\\
  &\leq32w\bigl(2^{d-2}c_{d-1}\Lambda^{d-2}+1\bigr)
   \leq c_dw\Lambda^{d-2},
\end{align*}
using \(\Lambda^{d-2}\geq1\).
\end{proof}
\end{annexitem}

\begin{annexitem}{annex:p03-w03-p03-064-boppana}{thm:p03-w03-p03-064-boppana}{Total influence of bounded-depth circuits (Boppana)}
\begin{proof}
Let \(C\) be a depth-\(d\) circuit of size \(s\geq2\) computing \(f\).  Form
\(C^+\): keep layer \(0\); insert a new layer \(1\) with one gate for each
literal node, whose only input is that literal and whose label is the one
not used on layer \(1\) of \(C\); move layer \(j\) of \(C\) to layer \(j+1\),
replacing each arc from a literal to a layer-\(1\) gate of \(C\) by an arc
from the new gate on that literal.  A one-input gate computes its input, so
\(C^+\) computes \(f\); the labels alternate; \(C^+\) has depth \(d+1\) and
width \(1\); and its layers \(2,\ldots,d\) are layers \(1,\ldots,d-1\) of
\(C\), with \(s\) gates.  By \cref{thm:p03-w03-p03-064-depth-induction} with
\(w=1\) and \(S=s\),
\(\mathbf I[f]\leq c_{d+1}(\log_2(4s))^{d-1}\).  For \(s\geq2\),
\(\log_2(4s)=2+\log_2s\leq3\log_2s\), so
\(\mathbf I[f]\leq3^{d-1}c_{d+1}(\log_2s)^{d-1}=K_d(\log_2s)^{d-1}\).
\end{proof}
\end{annexitem}
